\documentclass[11pt]{article}
\usepackage{amsmath}
\usepackage{hyperref}

\title{Untitled}
\author{Your Name}
\date{\today}

\begin{document}
\maketitle

\section{Introduction}\label{sec:intro}
This is a LaTeX document. Write \textbf{bold}, \emph{italic}, \texttt{code} and
\href{https://www.overleaf.com/learn}{links}.

\begin{itemize}
  \item A bullet list
  \item with two items
\end{itemize}

\begin{enumerate}
  \item A numbered list
  \item with two items
\end{enumerate}

\subsection{Math}
Inline math sits in dollar signs, like $a^2 + b^2 = c^2$. A numbered equation:
\begin{equation}\label{eq:sum}
  \sum_{k=1}^{n} k = \frac{n(n + 1)}{2}
\end{equation}

\subsection{Tables}
\begin{table}[h]
  \centering
  \begin{tabular}{lr}
    Item & Value \\
    \hline
    Alpha & 1 \\
    Beta & 2 \\
  \end{tabular}
  \caption{A small table.}\label{tab:items}
\end{table}

Table~\ref{tab:items} and Equation~\ref{eq:sum} are numbered for you.
Section~\ref{sec:intro} is the start.
\end{document}
